\documentclass[aspectratio=169,11pt]{beamer}
\usepackage[utf8]{inputenc}
\usepackage[T1]{fontenc}
\usepackage[french]{babel}
\usepackage{helvet}
\renewcommand{\familydefault}{\sfdefault}
\usepackage{hyperref}
\definecolor{bleu}{HTML}{0000FF}
\definecolor{fond}{HTML}{F4F4EF}
\setbeamercolor{background canvas}{bg=fond}
\setbeamercolor{structure}{fg=bleu}
\setbeamercolor{title}{fg=bleu}
\setbeamercolor{frametitle}{fg=bleu,bg=fond}
\setbeamercolor{normal text}{fg=black}
\setbeamertemplate{navigation symbols}{}
\setbeamertemplate{itemize item}{\textcolor{bleu}{\tiny$\blacksquare$}}
\setbeamertemplate{footline}{\hfill\footnotesize\insertframenumber\hspace{1em}\vskip6pt}
\hypersetup{colorlinks=true,urlcolor=bleu,linkcolor=bleu}
\title{Commissions économiques régionales de l'ONU}
\subtitle{Fonctionnement et données, vue d'ensemble}
\author{Veille institutionnelle}
\date{2026-11-03}

\begin{document}
\begin{frame}\titlepage\end{frame}

\section{Histoire et mission}
\begin{frame}{Histoire et mission}
\begin{itemize}
\item Cinq commissions sous l'autorité du Conseil économique et social des Nations unies
\item CEE (Genève) 1947, CESAP (Bangkok) 1947, CEPALC (Santiago) 1948, CEA (Addis-Abeba) 1958, CESAO (Beyrouth) 1973
\item Mission : coopération économique régionale et mise en œuvre des objectifs de développement durable
\item Fonctions : analyse, normes, forum intergouvernemental, assistance technique
\end{itemize}
\end{frame}

\begin{frame}{Histoire et mission : ressources à suivre}
\small
Veille hebdomadaire.
\begin{itemize}
\item \href{https://www.regionalcommissions.org}{Commissions régionales des Nations unies}\\{\footnotesize Site commun, publication : mensuelle}
\end{itemize}
\end{frame}

\section{Gouvernance}
\begin{frame}{Gouvernance}
\begin{itemize}
\item Chaque commission : États membres de la région, session annuelle ou biennale
\item Secrétaire exécutif nommé par le Secrétaire général de l'ONU
\item Rattachement au Conseil économique et social
\item Bureau de New York pour la coordination des commissions
\end{itemize}
\end{frame}

\begin{frame}{Gouvernance : ressources à suivre}
\small
Veille hebdomadaire.
\begin{itemize}
\item \href{https://www.regionalcommissions.org}{Site commun des commissions}\\{\footnotesize Site commun, publication : mensuelle}
\end{itemize}
\end{frame}

\section{Financement}
\begin{frame}{Financement}
\begin{itemize}
\item Budget ordinaire des Nations unies
\item Ressources extrabudgétaires de bailleurs
\item Compte pour le développement des Nations unies
\item Les montants diffèrent fortement selon les commissions (à vérifier)
\end{itemize}
\end{frame}

\begin{frame}{Financement : ressources à suivre}
\small
Veille hebdomadaire.
\begin{itemize}
\item \href{https://www.regionalcommissions.org}{Site commun des commissions}\\{\footnotesize Site commun, publication : trimestrielle}
\end{itemize}
\end{frame}

\section{Les cinq commissions}
\begin{frame}{Les cinq commissions (1/2)}
\begin{itemize}
\item CEA : Afrique, 54 pays membres (à vérifier)
\item CEE : Europe et Amérique du Nord, Asie centrale, 56 membres (à vérifier)
\item CEPALC : Amérique latine et Caraïbes
\item CESAP : Asie et Pacifique, 53 membres (à vérifier)
\end{itemize}
\end{frame}

\begin{frame}{Les cinq commissions (2/2)}
\begin{itemize}
\item CESAO : Asie occidentale, 20 pays arabes (à vérifier)
\end{itemize}
\end{frame}

\begin{frame}{Les cinq commissions : ressources à suivre}
\small
Veille hebdomadaire.
\begin{itemize}
\item \href{https://www.uneca.org}{CEA}\\{\footnotesize Site commission, publication : hebdomadaire}
\item \href{https://unece.org}{CEE}\\{\footnotesize Site commission, publication : hebdomadaire}
\item \href{https://www.cepal.org}{CEPALC}\\{\footnotesize Site commission, publication : hebdomadaire}
\item \href{https://www.unescap.org}{CESAP}\\{\footnotesize Site commission, publication : hebdomadaire}
\item \href{https://www.unescwa.org}{CESAO}\\{\footnotesize Site commission, publication : hebdomadaire}
\end{itemize}
\end{frame}

\subsection{CEA : Commission économique pour l'Afrique}
\begin{frame}{CEA : Commission économique pour l'Afrique}
\begin{itemize}
\item Siège à Addis-Abeba, créée en 1958
\item Travaille avec l'Union africaine et la BAD sur l'Agenda 2063 et la ZLECAf
\item Economic Report on Africa
\end{itemize}
\end{frame}

\begin{frame}{CEA : Commission économique pour l'Afrique : ressources à suivre}
\small
Veille hebdomadaire.
\begin{itemize}
\item \href{https://www.uneca.org}{CEA}\\{\footnotesize Site commission, publication : hebdomadaire}
\item \href{https://www.uneca.org/publications}{Publications de la CEA}\\{\footnotesize Catalogue, publication : mensuelle}
\end{itemize}
\end{frame}

\subsection{CEE : Commission économique pour l'Europe}
\begin{frame}{CEE : Commission économique pour l'Europe}
\begin{itemize}
\item Siège à Genève, créée en 1947
\item Normes et conventions : transports, environnement, commerce
\item Base statistique UNECE
\end{itemize}
\end{frame}

\begin{frame}{CEE : Commission économique pour l'Europe : ressources à suivre}
\small
Veille hebdomadaire.
\begin{itemize}
\item \href{https://unece.org}{CEE}\\{\footnotesize Site commission, publication : hebdomadaire}
\item \href{https://w3.unece.org/PXWeb2015/pxweb/en/STAT}{Base statistique UNECE}\\{\footnotesize Base de données, publication : mensuelle}
\end{itemize}
\end{frame}

\subsection{CEPALC : Amérique latine et Caraïbes}
\begin{frame}{CEPALC : Amérique latine et Caraïbes}
\begin{itemize}
\item Siège à Santiago, créée en 1948
\item CEPALSTAT : base statistique régionale
\item Étude économique annuelle de la région
\end{itemize}
\end{frame}

\begin{frame}{CEPALC : Amérique latine et Caraïbes : ressources à suivre}
\small
Veille hebdomadaire.
\begin{itemize}
\item \href{https://www.cepal.org}{CEPALC}\\{\footnotesize Site commission, publication : hebdomadaire}
\item \href{https://statistics.cepal.org}{CEPALSTAT}\\{\footnotesize Base de données, publication : mensuelle}
\end{itemize}
\end{frame}

\subsection{CESAP : Asie et Pacifique}
\begin{frame}{CESAP : Asie et Pacifique}
\begin{itemize}
\item Siège à Bangkok, créée en 1947
\item Portail de données statistiques de la CESAP
\item Étude économique et sociale annuelle
\end{itemize}
\end{frame}

\begin{frame}{CESAP : Asie et Pacifique : ressources à suivre}
\small
Veille hebdomadaire.
\begin{itemize}
\item \href{https://www.unescap.org}{CESAP}\\{\footnotesize Site commission, publication : hebdomadaire}
\item \href{https://data.unescap.org}{Données de la CESAP}\\{\footnotesize Base de données, publication : mensuelle}
\end{itemize}
\end{frame}

\subsection{CESAO : Asie occidentale}
\begin{frame}{CESAO : Asie occidentale}
\begin{itemize}
\item Siège à Beyrouth, créée en 1973 sous le nom de CEAO, renommée en 1985
\item Vingt pays arabes membres (à vérifier)
\end{itemize}
\end{frame}

\begin{frame}{CESAO : Asie occidentale : ressources à suivre}
\small
Veille hebdomadaire.
\begin{itemize}
\item \href{https://www.unescwa.org}{CESAO}\\{\footnotesize Site commission, publication : hebdomadaire}
\end{itemize}
\end{frame}

\section{Recherche et publications}
\begin{frame}{Recherche et publications}
\begin{itemize}
\item Economic Report on Africa (CEA)
\item Études économiques de la CEPALC : Estudio Económico de América Latina y el Caribe
\item Economic and Social Survey of Asia and the Pacific (CESAP)
\item Rapports sur les objectifs de développement durable par région
\end{itemize}
\end{frame}

\begin{frame}{Recherche et publications : ressources à suivre}
\small
Veille hebdomadaire.
\begin{itemize}
\item \href{https://www.uneca.org/publications}{Publications de la CEA}\\{\footnotesize Catalogue, publication : mensuelle}
\item \href{https://www.cepal.org/en/publications}{Publications de la CEPALC}\\{\footnotesize Catalogue, publication : mensuelle}
\end{itemize}
\end{frame}

\section{Données}
\begin{frame}{Données}
\begin{itemize}
\item Chaque commission tient ses propres bases statistiques régionales
\item UN Data et la base des indicateurs ODD de la Division de statistique
\item Méthodes comparables à vérifier avant tout croisement entre régions
\item Sous-dossiers : portails CEPALC, CESAP, CEE et ODD
\end{itemize}
\end{frame}

\begin{frame}{Données : ressources à suivre}
\small
Veille hebdomadaire.
\begin{itemize}
\item \href{https://data.un.org}{UN Data}\\{\footnotesize Portail, publication : mensuelle}
\item \href{https://unstats.un.org}{Division de statistique des Nations unies}\\{\footnotesize Portail, publication : mensuelle}
\item \href{https://unstats.un.org/sdgs}{Portail des données ODD}\\{\footnotesize Portail, publication : mensuelle}
\end{itemize}
\end{frame}

\section{Classifications}
\begin{frame}{Classifications}
\begin{itemize}
\item Découpage régional des Nations unies, norme M49
\item Régions géographiques différentes de celles des commissions
\item Pays membres de plusieurs commissions (États-Unis, Canada, pays d'Asie centrale)
\end{itemize}
\end{frame}

\begin{frame}{Classifications : ressources à suivre}
\small
Veille hebdomadaire.
\begin{itemize}
\item \href{https://unstats.un.org/unsd/methodology/m49/}{Norme M49}\\{\footnotesize Documentation, publication : annuelle}
\end{itemize}
\end{frame}

\section{Évaluation et contrôle}
\begin{frame}{Évaluation et contrôle}
\begin{itemize}
\item Bureau des services de contrôle interne de l'ONU
\item Corps commun d'inspection
\item Rapports au Conseil économique et social
\end{itemize}
\end{frame}

\begin{frame}{Évaluation et contrôle : ressources à suivre}
\small
Veille hebdomadaire.
\begin{itemize}
\item \href{https://www.unjiu.org}{Corps commun d'inspection}\\{\footnotesize Site de contrôle, publication : trimestrielle}
\end{itemize}
\end{frame}

\section{Plans stratégiques}
\begin{frame}{Plans stratégiques}
\begin{itemize}
\item Programme de développement durable à l'horizon 2030
\item Forums régionaux sur le développement durable chaque année
\item Cadres régionaux : Agenda 2063 pour l'Afrique, plans propres aux autres régions (à vérifier)
\end{itemize}
\end{frame}

\begin{frame}{Plans stratégiques : ressources à suivre}
\small
Veille hebdomadaire.
\begin{itemize}
\item \href{https://sdgs.un.org/2030agenda}{Programme 2030}\\{\footnotesize Page stratégique, publication : trimestrielle}
\end{itemize}
\end{frame}

\section{Critiques et débats}
\begin{frame}{Critiques et débats}
\begin{itemize}
\item Ressources limitées face aux mandats
\item Chevauchement avec d'autres institutions régionales
\item Dépendance aux financements extrabudgétaires
\end{itemize}
\end{frame}

\begin{frame}{Critiques et débats : ressources à suivre}
\small
Veille hebdomadaire.
\begin{itemize}
\item \href{https://www.regionalcommissions.org}{Site commun des commissions}\\{\footnotesize Site commun, publication : trimestrielle}
\end{itemize}
\end{frame}

\end{document}
